\documentclass[10pt,a4paper,roman]{moderncv}
\moderncvstyle{banking}
\moderncvcolor{purple}
\nopagenumbers{}
\usepackage[utf8]{inputenc}
\usepackage[T1]{fontenc}
\usepackage{lmodern}
\usepackage[scale=0.95,top=0.65cm,bottom=0.65cm,left=0.8cm,right=0.8cm]{geometry}
\setlength{\hintscolumnwidth}{2.25cm}
\setlength{\footskip}{18pt}
\name{Wahid}{SAMER}
\address{Cergy (95), France}{}{}
\mobile{07 53 10 95 74}
\email{samer.ahmed.wahid@gmail.com}
\social[linkedin]{wahid-samer-86b6b820a}
\social[github]{wahidrt}
\title{Alternance M2 - Market Risk / Quant Risk / Model Risk}
\begin{document}
\makecvtitle
\vspace{-3.2em}
\begin{center}
\parbox{0.97\textwidth}{\centering\small\itshape
Profil orienté \textbf{mesure des risques, modèles quantitatifs et produits dérivés}, avec une base solide en probabilités, simulation et méthodes numériques. Recherche d'une \textbf{alternance M2 de 12 mois à partir du 21 septembre 2026}. Mobilité : \textbf{France, Luxembourg, Suisse, Belgique}.}
\end{center}
\vspace{0.25em}

\section{Formation}
\cventry{2025 -- 2027}{Master Mathématiques Appliquées à l'Ingénierie Financière}{CY Cergy Paris Université / CY Tech}{Cergy}{}{\textbf{M1 :} probabilités, processus stochastiques, simulation stochastique, optimisation avancée, évaluation des actifs contingents et gestion de portefeuille.\\
\textbf{M2 2026--2027 :} grands risques et valeurs extrêmes, portfolio management, fixed income, model calibration and simulation, stochastic calculus.}
\cventry{2021 -- 2025}{Licence Mathématiques et Applications}{Université d'Évry Paris-Saclay}{Évry}{}{Formation solide en mathématiques pures et appliquées : probabilités, mesure et intégration, analyse, algèbre, topologie, analyse numérique et programmation.}

\section{Projets ciblés}
\cventry{2026}{Python}{Gestion du risque d'un système de trading}{Projet M1 - équipe de 4}{}{Suivi du portefeuille et du P\&L, exposition, coûts, drawdown, Sharpe/Sortino, turnover ; étude de contraintes de risque et de scénarios de marché.}
\cventry{2026}{Python / \LaTeX}{Black--Scholes, Heston et risque de modèle}{Mémoire M1}{}{Analyse des hypothèses de Black--Scholes, Greeks et sensibilité ; comparaison avec Heston pour intégrer une volatilité stochastique.}
\cventry{2025}{Excel / VBA}{Pricer d'options vanilles et stress tests}{Projet M1}{}{Pricing Call/Put, Greeks et scénarios de volatilité pour mesurer l'impact sur les prix et sensibilités.}

\section{Compétences}
\cvitem{Risque}{Mesures de performance/risque, stress tests, sensibilités (Greeks), exposition portefeuille, simulation.}
\cvitem{Mathématiques}{Probabilités, processus stochastiques, optimisation, valeurs extrêmes (M2), méthodes numériques.}
\cvitem{Programmation}{Python (Pandas, NumPy, SciPy, Scikit-learn, Matplotlib), VBA, SQL ; Excel, Git/GitHub.}
\cvitem{Langues}{Français : courant ; Arabe : langue maternelle ; Anglais : intermédiaire.}

\section{Expérience \& Engagement}
\cventry{Mai -- juin 2025}{Professeur de mathématiques - stagiaire}{Lycée Pierre Mendès France}{Vitrolles}{}{Préparation et transmission de notions de niveau Terminale ; développement de la rigueur, de la pédagogie et de la communication scientifique.}
\cventry{2019}{Technicien aéronautique - stagiaire}{G1 Aviation}{Gap}{}{Assemblage/maintenance dans un environnement exigeant en précision, procédures et sécurité.}
\cvitem{Associatif}{Bénévole aux \textbf{Restaurants du Cœur} pendant près de 3 ans : solidarité, accueil et travail d'équipe.}
\end{document}
